\begin{abstract}
Can self-supervised video models read the state of innate immune cells from label-free
video? We train TiJepa, a video joint-embedding predictive architecture, on
\val{pre.t1.clips} unlabelled clips of migrating human neutrophils, together with a
matched-step spatial-masking-only ablation, and evaluate frozen features against image
models, a frozen general-purpose video model (V-JEPA~2), random-initialisation and
frame-shuffle controls on every public dataset we found that pairs innate-cell imaging with
a molecular or functional readout, organised as a modality coverage map. Model choices use
only a label-free re-identification set held out from all tests. On microglia morphology TiJepa-1 is far above hand-crafted features and not
distinguishable from image models, while V-JEPA~2 exceeds it (balanced accuracy
\val{bb.E7_vjepa2_vitl.3c} vs \val{bb.E4_tijepa1.3c}); on cell re-identification it is
below raw pixels. Its predictor's apparent forecasting skill
is explained by positional structure in the latent space: it keeps \val{fc.E4_tijepa1.1.adv.static}
of its advantage over copying the last latent on clips in which nothing moves. Position
alone associates cells across frames better than any embedding; label-free autofluorescence
identity and NAD(P)H lifetime are predicted no better than by intensity statistics once
experimental condition is controlled. The representation is partially collapsed (pooled
effective rank \val{dg.E4_tijepa1.target.rank} of 384), and rank tracks downstream
performance across models. Public data pair innate-cell imaging with memory-relevant
readouts nowhere, so we close with a replicated transcript-level observation from
restimulation datasets and the design of the imaging study that would test memory
directly.
\end{abstract}
